%! Author = lili
%! Date = 8/20/26

\newglossaryentry{x_chromosom}{
    name={X-Chromosom},
    description={ Eines der beiden Geschlechtschromosomen bei Säugetieren, bei Frauen in zwei Kopien (XX), bei Männern in einer Kopie (XY) vorhanden.},
    type=defi
}

\newglossaryentry{uas}{
    name={UAS (Upstream Activating Sequence)},
    plural={UAS (Upstream Activating Sequences)},
    description={Regulatorische DNA-Sequenz stromaufwärts eines Promotors, an die Aktivatorproteine binden und dadurch die Transkription fördern. Besonders gut charakterisiert in Hefen.},
    type=defi
}

\newglossaryentry{urs}{
    name={URS (Upstream Repressing Sequence)},
    plural={URS (Upstream Repressing Sequences)},
    description={Regulatorische DNA-Sequenz stromaufwärts eines Promotors, an die Repressorproteine binden und dadurch die Transkription hemmen.},
    type=defi
}


\newglossaryentry{uce}{
    first={Upstream Control Element (UCE)},
    name={Upstream Control Element},
    description={Regulatorisches Promotorelement der RNA-Polymerase I, das stromaufwärts des Transkriptionsstarts liegt und die Transkription aktiviert},
    type=defi
}

\newglossaryentry{rrna}{
    name={rRNA},
    description={Ribosomale RNA; struktureller und katalytischer Bestandteil der Ribosomen, der maßgeblich an der Peptidbindung während der Translation beteiligt ist.},
    type=defi
}


\newglossaryentry{bre}{
    name={BRE (TFIIB Recognition Element)},
    plural={BRE (TFIIB Recognition Elements)},
    description={Core-Promotorelement, das von dem allgemeinen Transkriptionsfaktor TFIIB erkannt wird und die Bildung des Präinitiationskomplexes unterstützt.},
    type=defi
}

\newglossaryentry{pabp}{
    name={PABP (Poly(A)-binding protein)},
    plural={PABPs},
    description={Protein, das an den Poly(A)-Schwanz eukaryotischer mRNA bindet, diese stabilisiert, den Abbau hemmt und die Translationseffizienz erhöht.},
    type=defi
}
\newglossaryentry{peptidbindung}{
    name={Peptidbindung},
    plural={Peptidbindungen},
    description={Kovalente Bindung zwischen der Carboxylgruppe einer Aminosäure und der Aminogruppe einer anderen, die während der Translation im Ribosom gebildet wird.},
    type=defi
}

\newglossaryentry{tatabox}{
    name={TATA-Box},
    description={AT-reiche Promotorsequenz eukaryotischer Gene, meist etwa 25 Basenpaare vor dem Transkriptionsstart},
    type=defi
}

\newglossaryentry{tbp}{
    name={TBP (TATA-box binding protein)},
    description={Protein des TFIID-Komplexes, das an die TATA-Box bindet und die Bildung des Präinitiationskomplexes für die Transkription einleitet.},
    type=defi
}
